\section{Modelling decisions}\label{app:decisions}

Table~\ref{tab:decisions} records each point where the given data left a choice open, what we chose and why.

\begin{table}[!htbp]\centering\small
\begin{tabularx}{\textwidth}{>{\raggedright\arraybackslash}p{3.3cm}>{\raggedright\arraybackslash}p{4.2cm}X}
\toprule
Question & Decision & Reason\\
\midrule
What sets the pressure ratio? & $\Ph$ is fixed at the compressor outlet and $\Pl$ at its inlet. The ratio is a result. & The compressor is the only component that raises pressure, so its outlet is the highest pressure in the cycle. The \SI{20}{MPa} limit must apply there.\\
Recuperator: fix effectiveness or pinch? & Effectiveness 0.90. Reject designs with pinch below \SI{5}{K}. & 0.90 is given as the baseline and \SI{5}{K} as a minimum. Both are honoured without tuning each design.\\
How is the cooler sized? & Minimum \cotwo--water difference of \SI{5}{K}. & Only the water temperature is given. A pinch rule fixes the water flow, and the flow fixes $UA$, with nothing left to guess.\\
Cooler mass & \SI{2.5}{kg} per \si{kW/K}, as for recuperators. & No figure is given for coolers. The same exchanger technology is assumed.\\
Baseline mass split & $M_0$ calibrated at the most efficient feasible recuperated design. & \SI{60}{t} is given for the baseline layout. Calibrating once lets every other design be measured against it on the same basis.\\
Is the given 10~t recuperator figure used? & No. & It appears only inside a worked example. Exchanger masses are computed from $UA$ here.\\
Cost model & Cost $= \SI{250}{US\$/kg} \times$ mass. & A single price is given for installed hardware.\\
Economic measure & Net return over $N = 10$ years. & It puts a one-time cost and a yearly income on one scale, and needs no reference design.\\
Objectives & Lifetime payout and mass. & Payout combines efficiency and cost. Mass represents compactness. The critical-point margin is treated as a constraint and reported.\\
Carnot efficiency & Computed per design from $\Tin$ and $\Tt$. & Both temperatures are swept, so a single constant would misstate the second-law efficiency.\\
Turbine inlet temperature & Swept at and below \SI{500}{\celsius}. & It is given as an upper limit.\\
Is the simple cycle ranked? & No. It is solved and reported, but left out of the Pareto front, knee and TOPSIS. & At 15\% efficiency it is not a candidate, but as the lightest layout it would set the light end of the front and move the knee towards light designs.\\
What is layout 5? & Layout 3 with two intercooled main-compression stages. & The layout is given only as recompression with intercooling. The recompressor takes warm fluid, so intercooling is applied to the main compressor only.\\
\SI{7.7}{MPa} limit in layout 6 & Applied to the main compressor inlet only; the pre-compressor inlet ($\Pe$) is below it. & The layout requires $\Pe$ below the main compressor inlet pressure. The inlet stays single-phase, since its temperature is above the critical temperature.\\
The literal reading of that limit & Also ranked (strict ranking), but not followed. & It decides the result, so its outcome is reported beside ours.\\
When to stop refining & When the winners stop moving and none lies on a free edge of its range. & Further points in the same ranges would not change the answer.\\
$M_0$ in the studies & Recalibrated on the same reference design when the cycle model changes (S1); held at \SI{39.06}{t} where recalibration would switch the reference design (S3). & $M_0$ stands for the hardware outside the exchangers, which a change of reference design does not alter. A constant shift of every mass leaves the Pareto front and knee unchanged but moves TOPSIS, which scales by vector length.\\
\bottomrule
\end{tabularx}
\caption{Modelling decisions.}\label{tab:decisions}
\end{table}
